\documentclass{article}
\usepackage{hyperref}
\usepackage{listings}
\usepackage{xcolor}
\usepackage{enumitem}

\definecolor{codegreen}{rgb}{0,0.6,0}
\definecolor{codegray}{rgb}{0.5,0.5,0.5}
\definecolor{codepurple}{rgb}{0.58,0,0.82}
\definecolor{backcolour}{rgb}{0.95,0.95,0.92}

\lstdefinestyle{mystyle}{
    backgroundcolor=\color{backcolour},   
    commentstyle=\color{codegreen},
    keywordstyle=\color{magenta},
    numberstyle=\tiny\color{codegray},
    stringstyle=\color{codepurple},
    basicstyle=\ttfamily\footnotesize,
    breakatwhitespace=false,         
    breaklines=true,                 
    captionpos=b,                    
    keepspaces=true,                 
    numbers=left,                    
    numbersep=5pt,                  
    showspaces=false,                
    showstringspaces=false,
    showtabs=false,                  
    tabsize=2
}

\lstset{style=mystyle}

\title{FFXIV Character Viewer - Technical Documentation}
\author{Cascade AI}
\date{\today}

\begin{document}

\maketitle

\section{Overview}
The FFXIV Character Viewer is a web application that allows users to search for and view character information from Final Fantasy XIV's Lodestone website. This document provides a detailed explanation of how the application works, including example URLs and the data flow process.

\section{System Architecture}
The application consists of three main components:
\begin{enumerate}
    \item Frontend (React/JavaScript)
    \item Local Proxy Server (Node.js/Express)
    \item FFXIV Lodestone (External Data Source)
\end{enumerate}

\section{Search Process}

\subsection{Step 1: User Input}
When a user searches for a character (e.g., "Mommy Thames" on Aether), the following process occurs:

\begin{enumerate}[label=\alph*)]
    \item User enters the character name: "Mommy Thames"
    \item User selects or enters the Data Center: "Aether"
    \item The search form submits these parameters to our local API
\end{enumerate}

\subsection{Step 2: Local Proxy Server}
The server constructs a Lodestone URL with the following format:

\begin{lstlisting}[language=text]
https://na.finalfantasyxiv.com/lodestone/character/?q=Mommy%20Thames&worldname=_region_2&classjob=&race_tribe=&blog_lang=ja&blog_lang=en&blog_lang=de&blog_lang=fr&order=
\end{lstlisting}

Parameters explained:
\begin{itemize}
    \item \texttt{q=Mommy\%20Thames} - The character name (URL encoded)
    \item \texttt{worldname=\_region\_2} - North American data centers
    \item \texttt{blog\_lang} parameters - Include results in all languages
\end{itemize}

\subsection{Step 3: Request Headers}
The server adds browser-like headers to avoid being blocked:

\begin{lstlisting}[language=javascript]
headers = {
    'User-Agent': 'Mozilla/5.0 (Windows NT 10.0; Win64; x64) ...',
    'Accept': 'text/html,application/xhtml+xml,...',
    'Accept-Language': 'en-US,en;q=0.9',
    'Referer': 'https://na.finalfantasyxiv.com/lodestone/'
}
\end{lstlisting}

\section{Character Search Results}

\subsection{Step 1: HTML Parsing}
When the Lodestone responds, the application:

\begin{enumerate}[label=\alph*)]
    \item Receives the HTML content
    \item Uses DOMParser to parse the HTML
    \item Looks for character entries using the selector \texttt{div.entry}
    \item For each entry, extracts:
    \begin{itemize}
        \item Character ID from \texttt{a[href\^="/lodestone/character/"]}
        \item Name from \texttt{.entry\_\_name}
        \item World from \texttt{.entry\_\_world}
        \item Avatar from \texttt{.entry\_\_chara\_\_face img}
    \end{itemize}
\end{enumerate}

\section{Character Details Process}

\subsection{Step 1: Basic Information}
When a user selects a character, the application:

\begin{enumerate}[label=\alph*)]
    \item Fetches the character's profile page:
    \begin{lstlisting}[language=text]
https://na.finalfantasyxiv.com/lodestone/character/[ID]/
    \end{lstlisting}
    \item Extracts basic information using selectors:
    \begin{itemize}
        \item Name: \texttt{.frame\_\_chara\_\_name}
        \item Server: \texttt{.frame\_\_chara\_\_world}
        \item Portrait: \texttt{.frame\_\_chara\_\_face img}
        \item Biography: \texttt{.character\_\_selfintroduction}
    \end{itemize}
\end{enumerate}

\subsection{Step 2: Job Information}
The application then fetches detailed job information:

\begin{enumerate}[label=\alph*)]
    \item Requests the class/job page:
    \begin{lstlisting}[language=text]
https://na.finalfantasyxiv.com/lodestone/character/[ID]/class_job/
    \end{lstlisting}
    \item Parses the job data structure:
    \begin{itemize}
        \item Finds the main content container: \texttt{.character-content}
        \item Locates role sections: \texttt{.character\_job\_role}
        \item Gets role names from: \texttt{h4.heading--lead}
        \item For each job in \texttt{.character\_job.clearfix}:
        \begin{itemize}
            \item Name from \texttt{.character\_job\_name.js\_tooltip}
            \item Level from \texttt{.character\_job\_level}
            \item Experience from \texttt{.character\_job\_exp}
            \item Icon from \texttt{.character\_job\_icon img}
        \end{itemize}
    \end{itemize}
\end{enumerate}

\section{Error Handling}

\subsection{Server-Side}
The application includes robust error handling:

\begin{enumerate}
    \item Network request validation
    \begin{itemize}
        \item Checks response status codes
        \item Validates HTML content structure
        \item Returns appropriate error messages
    \end{itemize}
    
    \item Data parsing safeguards
    \begin{itemize}
        \item Validates HTML structure before parsing
        \item Provides fallback values for missing data
        \item Removes empty role sections
    \end{itemize}
    
    \item Response formatting
    \begin{itemize}
        \item Returns consistent error objects
        \item Includes detailed error messages
        \item Maintains valid response structure even on errors
    \end{itemize}
\end{enumerate}

\subsection{Client-Side}
The frontend handles errors gracefully:

\begin{enumerate}
    \item Search errors
    \begin{itemize}
        \item Displays user-friendly error messages
        \item Maintains UI state
        \item Allows retrying failed searches
    \end{itemize}
    
    \item Character details errors
    \begin{itemize}
        \item Shows fallback UI for missing data
        \item Handles partial data loading
        \item Provides feedback for failed requests
    \end{itemize}
\end{enumerate}

\section{Example URLs}

\subsection{Search URL}
For searching "Mommy Thames" on Aether:
\begin{lstlisting}[breaklines=true]
https://na.finalfantasyxiv.com/lodestone/character/?q=Mommy%20Thames&worldname=_region_2&classjob=&race_tribe=&blog_lang=ja&blog_lang=en&blog_lang=de&blog_lang=fr&order=
\end{lstlisting}

\subsection{Character Details URL}
For viewing a specific character (example ID: 12345678):
\begin{lstlisting}[breaklines=true]
https://na.finalfantasyxiv.com/lodestone/character/12345678/
\end{lstlisting}

\subsection{Job Details URL}
For viewing a character's job levels:
\begin{lstlisting}[breaklines=true]
https://na.finalfantasyxiv.com/lodestone/character/12345678/class_job/
\end{lstlisting}

\end{document}
